\section{Quick GIT guide and CSS conventions for the project}
\label{sec:coordination}

\subsection{How to use Git (part 2!): branches and pull requests}
\label{sec:git}

Our project lives in a Git repository hosted on GitHub
(\url{https://github.com/vittodevit/sa1}). This short guide explains how we
work with it: if you know how to clone, pull and push, you already know
enough to follow it.

\subsubsection{\texttt{main} and your other branches}
\label{sec:git-idea}

The repository has one special branch, called \texttt{main}. \texttt{main}
is the \textbf{real project}: the version of the website that is complete and
working, the one we will hand in at the end. Nobody ever works directly on
\texttt{main}.

Every time you want to modify something (a new page, a style fix, whatever),
you create a \textbf{branch}: a personal copy of the project that lives next
to \texttt{main} inside the same repository. While you work on your branch,
the rest of the group does not see what you do, and you cannot break the real
project. When your work is finished and checked, it is \textbf{merged} into
\texttt{main} through a \textbf{pull request}, and from that moment it is
part of the real project.

So every change always goes through the same life cycle:
\begin{itemize}
\item you take a fresh copy of the project by creating a branch,
\item you do your work there, saving it with commits,
\item you ask for your work to be merged back into \texttt{main} with a
pull request on GitHub,
\item after the merge the branch is not needed anymore, and your next task
starts from a fresh copy again.
\end{itemize}

\subsubsection{Step 1: get the latest version of the project}
\label{sec:git-step1}

If it is the first time you use the repository, clone it (this downloads a
full copy on your computer, and you only do it once):
\begin{verbatim}
git clone git@github.com:vittodevit/sa1.git
\end{verbatim}
The address is the one shown by the green \texttt{Code} button on the GitHub
page of the repository.\\

{\small Note: what you are seeing here is SSH cloning. For this to work you
need to have an SSH keypair configured on your Mac and GitHub account.
I suggest setting it up by searching a tutorial because it simplifies some
of the operations but token auth showed in the lectures works the same
(obviousely if you clone from HTTPS)}\\

If you already cloned it in the past, make sure you are on \texttt{main} and
update it, so that you start your work from the freshest version, the one
that includes what the others have merged in the meantime:
\begin{verbatim}
git checkout main
git pull
\end{verbatim}

\subsubsection{Step 2: create a branch and move to it}
\label{sec:git-step2}

Before touching any file, create a branch for the task you are going to work
on and switch to it. Pick a short, descriptive name, for example
\texttt{home-page} or \texttt{fix-footer}:
\begin{verbatim}
git branch <branch-name>
git checkout <branch-name>
\end{verbatim}
From this moment you are working on your copy: \texttt{main} stays safe, and
you can experiment freely.

If you are not sure on which branch you currently are, \texttt{git branch}
prints all the branches and marks the current one with a \texttt{*} (press 'Q' to exit).

\subsubsection{Step 3: do the work and commit it}
\label{sec:git-step3}

Now write your code (HTML, CSS, ...). Every time you reach a state you are
happy with, save a checkpoint with a commit:
\begin{verbatim}
git add .
git commit -m "short description of what you did"
\end{verbatim}
One commit per small working step is ideal: you will usually make one or
more commits before your task is done.

\subsubsection{Step 4: push your branch to GitHub}
\label{sec:git-step4}

When the work is done (or when you simply want to back it up), publish your
branch on GitHub with:
\begin{verbatim}
git push --set-upstream origin <branch-name>
\end{verbatim}
The first time you push a branch you need \texttt{--set-upstream}: it tells
Git to create a branch with the same name on GitHub and to connect yours to
it. For all the following pushes from the same branch, a simple
\texttt{git push} is enough.

\subsubsection{Step 5: open the pull request on GitHub}
\label{sec:git-step5}

Now go to the
\href{https://github.com/vittodevit/sa1}{repository page on GitHub}
and ask for your work to be merged:
\begin{enumerate}
\item open the \texttt{Pull requests} tab at the top of the page;
\item click on the green \texttt{New pull request} button;
\item with the two selectors at the top of the page, choose \texttt{main} as
the \emph{first} one (where your work should end up) and your branch as the
\emph{second} one (what you want to merge): the page shows a summary of the
changes you made;
\item click on \texttt{Create pull request}, write a meaningful title and
confirm.
\end{enumerate}

A pull request (PR) is simply the request ``I finished something on my
branch, please merge it into \texttt{main}''. It is also the place where the
work is checked before it becomes part of the real project.

\subsubsection{Step 6: if changes are requested, just push again}
\label{sec:git-step6}

Vittorio reviews every pull request. If there are issues, he will comment on
the PR and ask you to modify it. In that case \textbf{do not} close the PR
and \textbf{do not} create a new branch: your branch is still the right place
to work in. Just make sure you are on your branch, fix what was requested,
and commit and push as usual:
\begin{verbatim}
git checkout <branch-name>
git add .
git commit -m "fix the requested changes"
git push
\end{verbatim}
The pull request updates itself automatically with the new commits, and the
work continues on the same PR until it is accepted and merged.

\subsubsection{Step 7: after the merge}
\label{sec:git-step7}

When the PR is merged, your work is officially part of \texttt{main}. Before
starting anything new, go back to \texttt{main} and pull, so that your next
branch starts from a version that contains everybody's work, yours included:
\begin{verbatim}
git checkout main
git pull
\end{verbatim}
GitHub will offer to delete the merged branch: you can safely accept, it is
not needed anymore.

\subsubsection{If anything goes wrong}
\label{sec:git-troubles}

Git can be confusing at the beginning: a commit ended up on the wrong branch,
a command printed something scary, GitHub talks about
\emph{merge conflicts}, ... Do not try risky commands to ``fix'' things on your own:
it is much easier to fix these situations early.
Write your issue in the WhatsApp group of the project and
describe the problem, we will fix it together.

\subsection{Guide to the CSS}
\label{sec:css-guide}
\textit{TODO.}
